
\begin{table}[!htbp]
\centering
\caption{Seven-model pushover and performance-point summary.}
\label{tab:pushover_pp_summary_cp112}
\scriptsize
\resizebox{\textwidth}{!}{%
\begin{tabular}{@{}lcrrcrr@{}}
\toprule
Model & Peak dir. & Peak shear & Peak disp. & PP dir. & Max PP shear & Max PP disp. \\
 & & (kN) & (mm) & & (kN) & (mm) \\
\midrule
Regular & X & 1836.338 & 198.040 & X & 456.253 & 26.053 \\
L & Y & 1800.120 & 187.539 & X & 407.295 & 25.900 \\
L1 & X & 1606.250 & 212.287 & X & 340.753 & 24.592 \\
L2 & Y & 1804.152 & 172.789 & X & 370.935 & 21.171 \\
T & Y & 1640.715 & 223.937 & X & 343.488 & 24.802 \\
U & Y & 1882.070 & 174.697 & X & 435.084 & 27.072 \\
Non-Parallel & Y & 1827.694 & 154.530 & X & 399.636 & 21.328 \\
\bottomrule
\end{tabular}
}
\par\smallskip
\begin{minipage}{0.97\textwidth}\footnotesize
Table~\ref{tab:pushover_pp_summary_cp112} shows the peak pushover base shear and displacement for each building configuration, alongside the maximum reported performance-point values. The peak values describe the first global capacity peak on the recorded curves, whereas the performance-point quantities describe the available demand-state intersection. For all seven models, the reported maximum performance-point base shear and roof displacement occur in the X direction; these values should not be read as the corresponding peak-pushover ordinates.
\end{minipage}
\end{table}
